\documentclass[12pt,a4paper]{article}
\usepackage[utf8]{inputenc}
\usepackage[T1]{fontenc}
\usepackage{amsmath,amssymb,amsthm}
\usepackage{geometry}
\usepackage{enumitem}
\usepackage{hyperref}
\usepackage{fancyhdr}
\usepackage{titlesec}

\geometry{margin=1in}
\pagestyle{fancy}
\fancyhf{}
\rhead{Lecture 8}
\lhead{Quantum Systems via Linear Algebra}
\cfoot{\thepage}

\titleformat{\section}{\large\bfseries}{}{0em}{}
\titleformat{\subsection}{\normalsize\bfseries}{}{0em}{}

\newtheorem{definition}{Definition}
\newtheorem{theorem}{Theorem}

\title{\textbf{Lecture 8: Locality and the No-Signaling Theorem} \\ \large Unitarity, EPR, and Bell's Theorem}
\author{Quantum Systems via Linear Algebra, Differential Equations, and Probability}
\date{}

\begin{document}
\maketitle
\thispagestyle{fancy}

% ============================================================
\section{1. Reversibility of Entanglement}
% ============================================================

\begin{itemize}[nosep]
    \item Two subsystems that interact and become entangled \emph{can} be disentangled by reversing the interaction.  In principle, entanglement is \textbf{reversible} (unitary evolution is invertible).
    \item However, if one subsystem is lost to the environment, the remaining subsystem is stuck in a mixed state described by its density matrix.  From its perspective, the entanglement is effectively \textbf{irreversible}.
    \item Measurement is a special case: an apparatus with many degrees of freedom entangles with the system.  There is no practical way to reverse this---which is why measurements appear irreversible even though the underlying evolution is unitary.
\end{itemize}

% ============================================================
\section{2. The No-Signaling Theorem}
% ============================================================

\subsection{2.1 Statement}
No matter what Bob does to his subsystem---no matter what unitary operation he applies---Alice's density matrix is \textbf{unchanged}.

\begin{theorem}[No-signaling]
Let $v(a,b)$ be the component representation of the composite system.  If Bob applies a unitary matrix $U$ on his subsystem, the new components are
\[
    v'(a,b) = \sum_{b'} U_{bb'}\,v(a,b').
\]
Then Alice's density matrix satisfies
\[
    \rho'^{(A)}_{a,a'} = \rho^{(A)}_{a,a'}.
\]
\end{theorem}

\subsection{2.2 Proof}
Bob applies unitary $U$ to his subsystem.  The composite components become:
\[
    v'(a,b) = \sum_{b'} U_{bb'}\,v(a,b').
\]
Its complex conjugate (with index $a'$):
\[
    v'^*(a',b) = \sum_{\tilde{b}} U^*_{b\tilde{b}}\,v^*(a',\tilde{b}) = \sum_{\tilde{b}} U^\dagger_{\tilde{b}b}\,v^*(a',\tilde{b}).
\]
Now compute Alice's new density matrix:
\[
    \rho'^{(A)}_{a,a'} = \sum_b v'^*(a',b)\,v'(a,b)
    = \sum_b \sum_{\tilde{b}} \sum_{b'} U^\dagger_{\tilde{b}b}\,U_{bb'}\,v^*(a',\tilde{b})\,v(a,b').
\]
Perform the sum over $b$ first:
\[
    \sum_b U^\dagger_{\tilde{b}b}\,U_{bb'} = (U^\dagger U)_{\tilde{b}b'} = \delta_{\tilde{b}b'} \qquad (\text{unitarity!})
\]
This collapses the double sum, giving:
\[
    \rho'^{(A)}_{a,a'} = \sum_{b'} v^*(a',b')\,v(a,b') = \rho^{(A)}_{a,a'}.\;\qedsymbol
\]

\subsection{2.3 Physical interpretation}
\begin{itemize}[nosep]
    \item \textbf{Locality}: Bob cannot send any signal to Alice by manipulating his own degrees of freedom, regardless of entanglement.
    \item All of Alice's statistical predictions---expectation values, probability distributions---are completely unaffected by Bob's actions.
    \item This holds even for maximally entangled states.
    \item The proof \emph{critically depends on unitarity}.  If Bob's evolution were non-unitary ($U^\dagger U \neq I$), faster-than-light signaling would become possible---one deep reason why quantum mechanics resists modification.
\end{itemize}

% ============================================================
\section{3. The EPR Puzzle}
% ============================================================

Einstein, Podolsky, and Rosen (1935) highlighted the following:

\begin{quote}
A composite system can be in a state of complete knowledge (a pure state $\mathbf{v}$), while \emph{each} subsystem is in a state of complete ignorance (maximally mixed density matrix).  We know everything about the whole and nothing about the parts.
\end{quote}

\noindent Entanglement does \emph{not} allow signaling: the no-signaling theorem proves that Bob's manipulations leave Alice's density matrix unchanged.  The genuinely non-classical feature is that the correlations are stronger than any classical model can produce (Bell's theorem below).

% ============================================================
\section{4. Bell's Theorem}
% ============================================================

\subsection{4.1 The locality constraint}
Suppose Alice and Bob each have one subsystem of an entangled pair and are separated.  A classical model would assign pre-determined values to all measurement outcomes (``local hidden variables'').

\subsection{4.2 The theorem}
\textbf{Bell's theorem} shows that no assignment of pre-determined outcomes (no local hidden-variable theory) can reproduce all quantum correlations simultaneously for all measurement choices.

For a singlet state, if Alice measures $Z$ and gets $+1$, Bob must get $-1$.  But for \emph{all} possible measurement orientations simultaneously, the quantum correlations violate Bell inequalities---no classical pre-programming can match them.

This is the deep content: \textbf{no local hidden-variable theory can reproduce all predictions of quantum mechanics for entangled systems}.

% ============================================================
\section{5. Summary}
% ============================================================

\begin{enumerate}[nosep]
    \item Entanglement is in principle \textbf{reversible}, but practically irreversible when one subsystem is lost or has many degrees of freedom.
    \item The \textbf{no-signaling theorem}: Bob's unitary operations cannot change Alice's density matrix.  Locality is preserved.
    \item The proof depends on \textbf{unitarity}---modifying unitarity would allow faster-than-light signaling.
    \item The EPR puzzle: quantum ``knowledge of the whole'' does not imply ``knowledge of the parts,'' yet no signal can exploit this.
    \item \textbf{Bell's theorem}: no local hidden-variable theory can reproduce all quantum correlations for entangled systems.
\end{enumerate}

\end{document}
